% ----------------------------------------------------------------------
%  2.1.2. Энергия образования, концентрация и подвижность
% ----------------------------------------------------------------------
\subsection{Энергия образования точечных дефектов}\label{ssec:formation}

Этот пункт отвечает на три вопроса:
\begin{enumerate}[label=\arabic*)]
\item сколько энергии стоит создать дефект;
\item сколько дефектов есть в кристалле в равновесии;
\item как быстро дефекты перемещаются.
\end{enumerate}
Ответ на первый вопрос нужен для второго и третьего, поэтому начнём с него.

%=======================================================================
\subsubsection{Энергия образования вакансии}

\begin{opredelenie}
\term{Энергия образования точечного дефекта} $E^{\mathrm F}$~— работа,
которую нужно затратить, чтобы создать один дефект в идеальном кристалле.
Для вакансии ($\EvF$) подразумевается, что атом переносится из узла
внутри кристалла на его поверхность.
\end{opredelenie}

\paragraph{Простейшая оценка: разрыв связей}
Пусть атомы связаны только с ближайшими соседями, каждая связь имеет
энергию $\varepsilon$, а число ближайших соседей (координационное число)
равно $z$. Посчитаем энергию образования вакансии в два шага
(рис.~\ref{fig:schottky-count}).

\begin{shagi}
\item Вынимаем атом из объёма~— рвём $z$ связей. Затраты: $z\varepsilon$.
\item Ставим атом на поверхность, в «ступеньку» (излом), где у него
половина соседей. Выигрываем $z\varepsilon/2$.
\end{shagi}
Итого
\begin{equation}\label{eq:EvF-bond}
\EvF \approx z\varepsilon-\frac{z\varepsilon}{2}=\frac{z\varepsilon}{2}=E_c ,
\end{equation}
где $E_c=z\varepsilon/2$~— \term{энергия связи на атом}, или теплота
сублимации в расчёте на один атом: чтобы испарить кристалл, нужно разорвать
все связи, а каждая связь принадлежит двум атомам\ispr{В исходном тексте
(с.~155) из этих же соображений получена оценка $\EvF\approx 2H$. Она
соответствует удалению атома \emph{на бесконечность}, а не на поверхность
кристалла. Для вакансии, образованной по механизму Шоттки, оценка в модели
парных связей даёт $\EvF\approx H=E_c$, что и согласуется с выводом
на с.~152--153 оригинала.}.

\begin{figure}[ht]
\centering
\begin{tikzpicture}[scale=0.55]
  % кристалл
  \foreach \x in {0,...,9}{\foreach \y in {0,...,4}{
    \pgfmathtruncatemacro{\code}{10*\x+\y}
    \ifnum\code=32 \node[vac=4.6mm] at (\x,\y) {};
    \else \node[atom=4.6mm] at (\x,\y) {};\fi}}
  % поверхностная ступенька
  \foreach \x in {0,...,5}{\node[atom=4.6mm] at (\x,5) {};}
  \node[atomB=4.6mm] at (6,5) {};
  \draw[arr, fmRed, line width=0.9pt] (3,2) .. controls (2.5,6.5) and (5.5,7) .. (5.85,5.4);
  \node[slbl, align=center, fmRed] at (2.2,6.6) {атом уходит\\на поверхность};
  \draw[decorate, decoration={brace, amplitude=4pt, mirror}] (9.6,-0.4) -- node[right=5pt, slbl, align=left]{объём:\\$z$ соседей} (9.6,4.4);
  \node[slbl, align=left, right] at (6.5,5.2) {излом ступеньки:\\$z/2$ соседей};
\end{tikzpicture}
\caption{Образование вакансии по механизму Шоттки: атом из объёма
кристалла переходит на поверхность, в положение излома ступеньки}
\label{fig:schottky-count}
\end{figure}

\paragraph{Почему простая оценка завышена}
Опыт даёт $\EvF$ в 2,5--3 раза меньше, чем $E_c$ (табл.~\ref{tab:EvF}).
Мы не учли два эффекта, которые \emph{снижают} энергию:
\begin{itemize}
\item \vazhno{релаксацию}: соседи вакансии сдвигаются к пустому узлу и
частично «залечивают» разорванные связи;
\item \vazhno{перераспределение электронов}: в металле электронный газ
заполняет пустоту и экранирует её.
\end{itemize}
Учитывая их, пишут $\EvF=E_c-W$, где $W$~— суммарная энергия релаксации.
Главная трудность теории~— именно расчёт $W$. Такие расчёты выполнили
Хантингтон, Зейтц, Фуми, Зегер и др.; их результаты для меди лежат
в пределах $0{,}8$--$1{,}2\eV$.

\begin{table}[ht]
\centering
\caption{Энергия связи $E_c$ и измеренная энергия образования вакансии
$\EvF$ некоторых металлов (справочные данные)}
\label{tab:EvF}
\small
\begin{tabular}{@{}lcccccccc@{}}
\toprule
 & Al & Ag & Au & Cu & Pt & Ni & Mo & W \\
\midrule
Решётка & ГЦК & ГЦК & ГЦК & ГЦК & ГЦК & ГЦК & ОЦК & ОЦК \\
$E_c$, эВ & 3,39 & 2,95 & 3,81 & 3,49 & 5,84 & 4,44 & 6,82 & 8,90 \\
$\EvF$, эВ & 0,67 & 1,11 & 0,93 & 1,28 & 1,35 & 1,79 & 3,0 & 3,7 \\
$\EvF/E_c$ & 0,20 & 0,38 & 0,24 & 0,37 & 0,23 & 0,40 & 0,44 & 0,42 \\
$\Tpl$, К & 933 & 1235 & 1337 & 1358 & 2041 & 1728 & 2896 & 3695 \\
$\EvF/k\Tpl$ & 8,3 & 10,4 & 8,1 & 10,9 & 7,7 & 12,0 & 12,0 & 11,6 \\
\bottomrule
\end{tabular}
\end{table}

\paragraph{Полезное эмпирическое правило}
Из табл.~\ref{tab:EvF} видно, что энергия образования вакансии
пропорциональна температуре плавления:
\begin{keyeq}\label{eq:EvF-Tm}
\EvF\approx (8\text{--}12)\,k\Tpl \approx 10\,k\Tpl .
\end{keyeq}
Это естественно: и плавление, и образование вакансии требуют
«расшатывания» межатомных связей, и оба определяются их прочностью.

\begin{podrobnee}[Почему энергия вакансии сосредоточена у самого дефекта]
За пределами нескольких межатомных расстояний решётку можно считать
сплошной упругой средой. Вакансию тогда заменяют \emph{центром дилатации}:
изнутри маленькой сферической полости радиусом $\sim b$ ($b$~— межатомное
расстояние) действует давление, которое создаёт на её поверхности
нужное смещение соседей. В таком поле смещения $u\propto r^{-2}$,
деформации и напряжения $\propto r^{-3}$, а плотность упругой энергии
$\propto r^{-6}$. Энергия поля за пределами сферы радиусом $R$
\[
U(>R)\propto\int_R^{\infty} r^{-6}\,r^{2}\,\d r\propto R^{-3}.
\]
Значит, вне сферы $R=5b$ находится лишь $(1/5)^3<1\,\%$ всей энергии.
Почти вся энергия вакансии сосредоточена в ближайшем окружении
($r\lesssim 2b$), где всего несколько десятков атомов. Поэтому $\EvF$
определяется дискретным атомным строением ядра дефекта, а не упругостью
сплошной среды.
\end{podrobnee}

%=======================================================================
\subsubsection{Энергия образования межузельного атома}

Межузельный атом приходится «втискивать» между соседями, поэтому решётка
вокруг него искажена гораздо сильнее, чем вокруг вакансии. Отсюда
\begin{glavnoe}
Энергия образования межузельного атома в несколько раз больше, чем
вакансии: $\EiF\approx 3$--$5\eV$ против $\EvF\approx 1\eV$ (для меди
$\EiF\approx 3\eV$).
\end{glavnoe}
Точное значение $\EiF$ известно хуже, чем $\EvF$: прямых равновесных
измерений нет, а расчёты чувствительны к выбору межатомного потенциала.

%=======================================================================
\subsubsection{Искажения решётки вокруг дефекта}

Когда в решётке появляется дефект, соседние атомы смещаются в положения
с наименьшей энергией. Это смещение называют \term{атомной релаксацией}.
Именно она заметно снижает энергии образования и миграции дефектов.

Точный расчёт релаксации непрост: сдвинув первую координационную сферу,
мы меняем силы, действующие на вторую, та~— на третью, и т.\,д.;
положения всех атомов приходится подбирать одновременно, минимизируя
энергию. Результаты расчётов (табл.~\ref{tab:relax}) показывают
следующее.
\begin{itemize}
\item Ближайшие соседи вакансии смещаются \emph{к} ней примерно на 2\,\%
межатомного расстояния (в ГЦК-меди). Соседи межузельного атома
сдвигаются \emph{от} него сильнее~— на 10--15\,\%.
\item Смещения убывают с расстоянием \vazhno{немонотонно}: во второй
координационной сфере вакансии они малы и имеют \emph{противоположный}
знак. Поле смещений вблизи дефекта анизотропно.
\item В ОЦК-металлах релаксация сильнее, чем в ГЦК: ОЦК-решётка не
плотноупакованная, и атомам «есть куда сдвинуться».
\end{itemize}
Область вокруг дефекта, где смещения велики, называют \term{ядром дефекта}.

\begin{table}[ht]
\centering
\caption{Расчётные смещения атомов первых трёх координационных сфер
вокруг точечных дефектов (в долях расстояния между ближайшими атомами;
знак «$-$»~— смещение к дефекту)}
\label{tab:relax}
\small
\begin{tabular}{@{}llccc@{}}
\toprule
Металл & Дефект & 1-я сфера & 2-я сфера & 3-я сфера \\
\midrule
Cu (ГЦК) & вакансия & $-0{,}021$ & $+0{,}002$ & $-0{,}002$ \\
Cu (ГЦК) & межузельный атом (окт.) & $+0{,}149$ & --- & $+0{,}032$ \\
Fe (ОЦК) & вакансия & $-0{,}061$ & $+0{,}021$ & $-0{,}003$ \\
Na (ОЦК) & межузельный атом & $+0{,}3$ & --- & --- \\
\bottomrule
\end{tabular}
\end{table}

%=======================================================================
\subsubsection{Равновесная концентрация точечных дефектов}\label{sssec:eq-conc}

Создание дефекта требует энергии~— казалось бы, в равновесии дефектов
быть не должно. Однако при $T>0$ равновесию отвечает не минимум энергии,
а минимум \emph{свободной энергии} $G=H-TS$. Дефекты увеличивают энтропию
(беспорядок), и при конечной температуре некоторое их количество
\emph{выгодно}.

\begin{glavnoe}
Равновесная концентрация дефектов~— результат компромисса: каждый дефект
стоит энергии $E^{\mathrm F}$ (это повышает $G$), но увеличивает число
способов расположения атомов, т.\,е. энтропию (это понижает $G$ на
$-TS$). Минимум $G$ достигается при конечной концентрации дефектов,
растущей с температурой (рис.~\ref{fig:G-of-n}).
\end{glavnoe}

\begin{figure}[ht]
\centering
\begin{tikzpicture}
\begin{axis}[width=0.62\textwidth, height=5.6cm,
  xmin=0, xmax=0.16, ymin=-0.5, ymax=0.5,
  axis lines=middle, xtick=\empty, ytick=\empty,
  xlabel={$n$}, ylabel={$\Delta G$},
  every axis x label/.style={at={(ticklabel* cs:1)}, anchor=west},
  every axis y label/.style={at={(ticklabel* cs:1)}, anchor=south},
  clip=false]
  \addplot[fmRed, thick, domain=0:0.16, samples=60] {3*x};
  \addplot[fmBlue, thick, domain=0.0005:0.16, samples=120] {x*ln(x)+(1-x)*ln(1-x)};
  \addplot[black, very thick, domain=0.0005:0.16, samples=120] {3*x+x*ln(x)+(1-x)*ln(1-x)};
  \draw[dashed, fmGray] (axis cs:0.0498,0) -- (axis cs:0.0498,-0.0486);
  \node[below, font=\small] at (axis cs:0.0498,-0.06) {$n_{\text{равн}}$};
  \node[fmRed, font=\small, anchor=west] at (axis cs:0.11,0.40) {$nE^{\mathrm F}$};
  \node[fmBlue, font=\small, anchor=west] at (axis cs:0.115,-0.33) {$-TS_{\text{конф}}$};
  \node[font=\small, anchor=west] at (axis cs:0.13,0.075) {$\Delta G$};
\end{axis}
\end{tikzpicture}
\caption{Изменение свободной энергии кристалла $\Delta G$ при образовании
$n$ дефектов (схема). Энергия растёт линейно ($nE^{\mathrm F}$), а
энтропийный вклад $-TS$ сначала падает очень круто. Сумма имеет минимум
при $n=n_{\text{равн}}$}
\label{fig:G-of-n}
\end{figure}

\paragraph{Конфигурационная энтропия}
По формуле Больцмана $S=k\ln W$, где $W$~— число способов, которыми можно
осуществить данное состояние. Атомы одного сорта неразличимы, поэтому
идеальный кристалл ($N$ атомов в $N$ узлах) можно построить единственным
способом: $W=1$, $S=0$.

Добавим вакансии и посчитаем $W$.
\begin{itemize}
\item \emph{Одна вакансия} может находиться в любом из $N$ узлов: $W_1=N$.
\item \emph{Две вакансии}: первую ставим $N$ способами, вторую $(N-1)$,
но вакансии неразличимы, и перестановка их местами не даёт новой
конфигурации; поэтому $W_2=N(N-1)/2$.
\item \emph{$n$ вакансий}: число способов выбрать $n$ узлов из $N$,
\begin{equation}\label{eq:Wn}
W_n=C_N^n=\frac{N!}{n!\,(N-n)!},\qquad S=k\ln\frac{N!}{n!\,(N-n)!}.
\end{equation}
\end{itemize}

\paragraph{Минимизация свободной энергии}
Изменение свободной энергии кристалла при образовании $n$ вакансий
\begin{equation}\label{eq:dG}
\Delta G = n\EvF - T\left(S+nS_v\right).
\end{equation}
Здесь $S$~— конфигурационная энтропия~\eqref{eq:Wn}, а $S_v$~— добавочная
\emph{колебательная} энтропия на одну вакансию: около вакансии связи
ослаблены, соседние атомы колеблются с меньшей частотой, и это тоже
увеличивает беспорядок. Концентрацию считаем малой, $n\ll N$, так что
дефекты не взаимодействуют друг с другом.

\begin{shagi}
\item \emph{Формула Стирлинга.} Для больших $X$: $\ln X!\approx X\ln X-X$.
Тогда
\[
S=k\bigl[N\ln N-n\ln n-(N-n)\ln(N-n)\bigr].
\]
\item \emph{Условие минимума} $\partial\Delta G/\partial n=0$:
\[
\EvF-TS_v-T\frac{\partial S}{\partial n}=0,\qquad
\frac{\partial S}{\partial n}=k\ln\frac{N-n}{n}.
\]
\item \emph{Решение:}
\[
\ln\frac{n}{N-n}=\frac{S_v}{k}-\frac{\EvF}{kT}
\quad\Longrightarrow\quad
\frac{n}{N-n}=\Ex{\frac{S_v}{k}}\Ex{-\frac{\EvF}{kT}}.
\]
\item Поскольку $n\ll N$, в знаменателе $N-n\approx N$.
\end{shagi}

В результате получаем главную формулу раздела:
\begin{keyeq}\label{eq:Cv}
C_v=\frac{n_v}{N}=\Ex{\frac{S_v}{k}}\Ex{-\frac{\EvF}{kT}}\approx\Ex{-\frac{\EvF}{kT}} .
\end{keyeq}
Колебательная энтропия вакансии $S_v\approx(1\text{--}2)\,k$, так что
множитель $\exp(S_v/k)\approx 3$--$7$. Он не меняет порядка величины, и
в оценках его часто опускают.

\begin{smysl}
$\exp(-\EvF/kT)$~— это \emph{больцмановский фактор}, вероятность
тепловой флуктуации с энергией $\EvF$. Формула~\eqref{eq:Cv} говорит:
доля пустых узлов равна вероятности того, что тепловое движение
«соберёт» в одном месте энергию, нужную для образования вакансии.
\end{smysl}

Та же формула справедлива и для межузельных атомов, если заменить
$\EvF$ на $\EiF$.

\paragraph{Порядки величин}
Концентрация дефектов зависит от температуры экспоненциально
(рис.~\ref{fig:Cv-T}).

\begin{figure}[ht]
\centering
\begin{tikzpicture}
\begin{axis}[width=0.66\textwidth, height=6cm,
  ymode=log, xmin=250, xmax=1400, ymin=1e-20, ymax=1e-3,
  xlabel={$T$, К}, ylabel={$C_v$},
  ytick={1e-20,1e-16,1e-12,1e-8,1e-4},
  grid=major, grid style={fmGray!20}]
  \addplot[fmBlue, very thick, domain=280:1380, samples=80] {exp(-1.1/(8.617e-5*x))};
  \addplot[only marks, mark=*, mark size=2pt, fmRed] coordinates
    {(300,3.3e-19) (700,1.2e-8) (1100,9.0e-6) (1350,7.8e-5)};
  \node[font=\footnotesize, anchor=west] at (axis cs:300,3.3e-19) {$\;3\cdot10^{-19}$};
  \node[font=\footnotesize, anchor=north west] at (axis cs:700,1.2e-8) {$10^{-8}$};
  \node[font=\footnotesize, anchor=north west] at (axis cs:1100,9e-6) {$10^{-5}$};
  \node[font=\footnotesize, anchor=south east] at (axis cs:1350,7.8e-5) {$10^{-4}$};
\end{axis}
\end{tikzpicture}
\caption{Равновесная концентрация вакансий при $\EvF=1{,}1\eV$ (близко к
значениям для многих металлов). От комнатной температуры до точки
плавления $C_v$ вырастает на 15 порядков}
\label{fig:Cv-T}
\end{figure}

\begin{primer}[Сколько вакансий в меди]
Для меди $\EvF\approx 1{,}28\eV$, $\Tpl=1358$~К.
\begin{itemize}
\item При $T=300$~К: $kT=0{,}0259\eV$, $\EvF/kT\approx 49$,
$C_v\approx e^{-49}\approx 4\cdot10^{-22}$. В кубическом сантиметре меди
($8{,}5\cdot10^{22}$ атомов)~— всего несколько равновесных вакансий.
\item При $T=\Tpl$: $kT=0{,}117\eV$, $\EvF/kT\approx 10{,}9$,
$C_v\approx e^{-10{,}9}\approx 2\cdot10^{-5}$; с учётом
$\exp(S_v/k)$~— около $10^{-4}$, т.\,е. одна вакансия на
$\sim 10^4$ узлов.
\end{itemize}
\end{primer}

\begin{glavnoe}
Вблизи температуры плавления концентрация вакансий примерно одинакова
для всех металлов и составляет $\sim 10^{-4}$ (это прямо следует из
правила $\EvF\approx 10\,k\Tpl$). Такая концентрация уже существенна:
например, она определяет скорость диффузии.
\end{glavnoe}

Для межузельных атомов $\EiF$ в несколько раз больше, и даже у точки
плавления их равновесная концентрация ничтожна: для меди
$\exp(-3/0{,}117)\approx 10^{-11}$. Поэтому \vazhno{в тепловом равновесии
межузельными атомами обычно можно пренебречь}, и равновесные точечные
дефекты в металле~— это вакансии.

%=======================================================================
\subsubsection{Вакансии в сплаве}

Атом замещения отличается от атомов основы размером и электронным
строением, поэтому искажает решётку. Образовать вакансию рядом с ним
легче или труднее, чем в чистом металле. Описывают это
\term{энергией связи} вакансии с атомом примеси $B_{iv}$: энергия
образования вакансии в соседнем с примесью узле равна $\EvF-B_{iv}$.
При $B_{iv}>0$ вакансия и примесь притягиваются. Например, если атом
замещения \emph{крупнее} атома основы, соседний пустой узел частично
снимает созданное им сжатие, и $B_{iv}>0$.

Пусть доля примесных атомов $C_i$ мала, а $z$~— координационное число.
Тогда доля узлов, соседних с примесными атомами, равна $zC_i$, а
остальных $(1-zC_i)$. Повторяя вывод~\eqref{eq:Cv} для каждой группы
узлов, получаем\ispr{В формуле (2.4) оригинала во втором слагаемом
пропущена экспонента.}
\begin{equation}\label{eq:Cv-alloy}
C_v=(1-zC_i)\Ex{-\frac{\EvF}{kT}}+zC_i\,\Ex{-\frac{\EvF-B_{iv}}{kT}} .
\end{equation}
Формула верна, пока первые координационные сферы разных примесных атомов
не перекрываются, т.\,е. при $zC_i\ll 1$.

Большинство вакансий окажется \emph{рядом с примесью}, если второе
слагаемое больше первого. При $zC_i\ll1$ это условие принимает вид
\begin{equation}\label{eq:Biv-cond}
\frac{B_{iv}}{kT}>\ln\frac{1}{zC_i}.
\end{equation}

\begin{primer}[Где сидят вакансии в разбавленном сплаве]
Пусть в ГЦК-металле ($z=12$) растворён 1\,ат.\,\% примеси: $C_i=0{,}01$,
$zC_i=0{,}12$, $\ln(1/0{,}12)\approx 2{,}1$. При $T=600$~К ($kT=0{,}052\eV$)
условие~\eqref{eq:Biv-cond} выполняется уже при $B_{iv}>0{,}11\eV$~— это
вполне реальная энергия связи. Значит, даже при 1\,\% примеси б\'{о}льшая
часть вакансий может находиться рядом с примесными атомами.
\end{primer}

%=======================================================================
\subsubsection{Миграция вакансии}

Вакансия перемещается так: один из её ближайших соседей перескакивает
в пустой узел. Атом сдвигается вперёд, вакансия~— назад, на его место
(рис.~\ref{fig:vac-jump}).

\begin{figure}[ht]
\centering
\begin{subfigure}[b]{0.42\textwidth}\centering
\begin{tikzpicture}[scale=0.8]
  \foreach \x in {0,...,4}{\foreach \y in {0,...,3}{
    \pgfmathtruncatemacro{\code}{10*\x+\y}
    \ifnum\code=21 \node[vac=5.2mm] at (\x,\y) {};
    \else\ifnum\code=31 \node[atomB=5.2mm] at (\x,\y) {};
    \else \node[atom=5.2mm] at (\x,\y) {};\fi\fi}}
  \draw[arr, fmRed, line width=1pt] (2.75,1.05) -- (2.25,1.05);
  \draw[arr, fmBlue, dashed] (2.25,0.6) -- (2.75,0.6);
  \node[slbl, fmBlue] at (2.5,0.35) {вакансия};
  \node[slbl, fmRed] at (2.5,1.35) {атом};
\end{tikzpicture}
\caption{}
\end{subfigure}
\begin{subfigure}[b]{0.42\textwidth}\centering
\begin{tikzpicture}[scale=0.8]
  \foreach \x in {0,...,4}{\foreach \y in {0,...,3}{
    \pgfmathtruncatemacro{\code}{10*\x+\y}
    \ifnum\code=21 \else\ifnum\code=31 \else\ifnum\code=32 \else\ifnum\code=30 \else
      \node[atom=5.2mm] at (\x,\y) {};\fi\fi\fi\fi}}
  \node[atom=5.2mm] at (3,2.17) {};
  \node[atom=5.2mm] at (3,-0.17) {};
  \node[vac=5.2mm] at (2,1) {};
  \node[vac=5.2mm] at (3,1) {};
  \node[atomB=5.2mm] at (2.5,1) {};
  \draw[arr, fmGray] (3,2.0) ++(0,0.45) -- ++(0,0.25);
  \draw[arr, fmGray] (3,-0.0) ++(0,-0.45) -- ++(0,-0.25);
\end{tikzpicture}
\caption{}
\end{subfigure}
\caption{Перескок атома в вакансию: \textit{а}~— исходное положение;
\textit{б}~— седловая точка: перескакивающий атом находится посередине
между узлами и должен «раздвинуть» атомы, запирающие проход (стрелки).
Именно в этот момент энергия максимальна}
\label{fig:vac-jump}
\end{figure}

Если построить энергию кристалла в зависимости от положения
перескакивающего атома (или, что то же, от положения вакансии), получится
периодическая кривая с периодом $a$ (рис.~\ref{fig:vac-profile}).
Минимумы соответствуют узлам решётки, максимум~— середине пути, где
искажения наибольшие. Высота барьера $\Evm$ обычно намного больше $kT$,
поэтому перескок происходит лишь тогда, когда тепловая флуктуация сообщит
атому нужную энергию.

\begin{opredelenie}
\term{Энергия активации миграции} вакансии $\Evm$~— высота
энергетического барьера, который нужно преодолеть при перескоке
вакансии в соседний узел.
\end{opredelenie}

\begin{figure}[ht]
\centering
\begin{tikzpicture}
\begin{axis}[width=0.6\textwidth, height=4.6cm,
  xmin=-0.1, xmax=3.1, ymin=0, ymax=1.75, axis lines=left,
  xtick={0,1,2,3}, xticklabels={$0$,$a$,$2a$,$3a$},
  ytick={0.7,1.5}, yticklabels={$\EvF$,{}},
  xlabel={положение вакансии $x$}, ylabel={$E$}]
  \addplot[fmBlue, very thick, domain=-0.05:3.05, samples=150]
    {0.7+0.4*(1-cos(deg(2*pi*x)))};
  \draw[dimline] (axis cs:1.5,0.7) -- node[right, font=\small]{$\Evm$} (axis cs:1.5,1.5);
  \draw[dashed, fmGray] (axis cs:-0.1,0.7) -- (axis cs:3.1,0.7);
  \draw[dashed, fmGray] (axis cs:1.2,1.5) -- (axis cs:1.8,1.5);
\end{axis}
\end{tikzpicture}
\caption{Энергия кристалла с вакансией в зависимости от её положения.
Узлы решётки~— минимумы, седловые точки посередине~— максимумы}
\label{fig:vac-profile}
\end{figure}

\paragraph{Частота перескоков}
Каждый атом колеблется около узла с частотой $\nu_0\approx
10^{12}$--$10^{13}$~с$^{-1}$ (порядка дебаевской). При каждом колебании
он «пробует» перескочить, и попытка успешна с вероятностью
$\exp(-\Evm/kT)$. Перескочить в вакансию может любой из её $z$ соседей,
поэтому частота перескоков вакансии
\begin{keyeq}\label{eq:Gamma-v}
\Gamma_v=z\nu_0\Ex{-\frac{\Evm}{kT}} .
\end{keyeq}
За каждый перескок вакансия смещается на межатомное расстояние $a$, и
скорость её блуждания по порядку величины равна $a\Gamma_v$.

\begin{primer}[Как часто прыгает вакансия в меди]
Примем $\Evm=0{,}7\eV$, $z=12$, $\nu_0=10^{13}$~с$^{-1}$.
\begin{itemize}
\item $T=300$~К: $\Evm/kT=27$, $\Gamma_v\approx 1{,}2\cdot10^{14}\cdot e^{-27}
\approx 200$~с$^{-1}$.
\item $T=1000$~К: $\Evm/kT=8{,}1$, $\Gamma_v\approx 1{,}2\cdot10^{14}\cdot
e^{-8{,}1}\approx 4\cdot10^{10}$~с$^{-1}$.
\end{itemize}
Нагрев с комнатной температуры до 1000~К ускоряет миграцию в
$\sim10^8$ раз.
\end{primer}

\paragraph{Сколько стоит миграция}
Из опыта для ГЦК-металлов $\Evm$ того же порядка, что и $\EvF$, но
несколько меньше (для меди $\Evm\approx 0{,}7$--$1{,}0\eV$). В ОЦК-металлах
энергия миграции вакансий меньше, чем в ГЦК: решётка более рыхлая,
межатомные расстояния больше, и при перескоке атому приходится меньше
«расталкивать» соседей. Поэтому \vazhno{в ОЦК-металлах вакансии более
подвижны}.

\begin{podrobnee}[Оценка $\Evm$ из кинетики перескока]
Пока атом проходит седловую точку (время порядка периода колебаний
$t\approx\nu_0^{-1}$), возмущение успевает распространиться со скоростью
звука лишь на $\sim 2a$. Значит, энергия $\Evm$ в основном кинетическая и
принадлежит самому перескакивающему атому: чтобы проскочить проход за
время $t$, пока соседи раздвинуты, атом массы $m$ должен приобрести
скорость $\upsilon\approx qa/t$. Тогда
\[
\Evm\approx\frac{m\upsilon^2}{2}=\frac{m}{2}\left(\frac{qa}{t}\right)^2 ,
\]
где множитель $q\approx 1$ учитывает, что скорость непостоянна, а проход
открыт не весь период. Лучшее согласие с опытом получается при
$q\approx 0{,}9$ для плотноупакованных решёток и $q\approx 0{,}8$ для ОЦК.
\end{podrobnee}

%=======================================================================
\subsubsection{Самодиффузия по вакансионному механизму}

Блуждание вакансий означает, что перемещаются и атомы: атом долго
колеблется около своего узла, затем к нему подходит вакансия, атом
перескакивает на её место и оказывается среди новых соседей.
Перемещение атомов по собственной решётке называется
\term{самодиффузией}. Механизмов самодиффузии можно предложить несколько,
но в металлах основной~— \term{вакансионный}.

\paragraph{Частота перескоков атома}
Чтобы данный атом перескочил, нужны два независимых события:
\begin{enumerate}[label=\arabic*)]
\item рядом с атомом должна оказаться вакансия; вероятность этого
$zC_v=z\exp(-\EvF/kT)$;
\item атом должен преодолеть барьер; частота успешных попыток
$\nu_0\exp(-\Evm/kT)$.
\end{enumerate}
Перемножая, получаем частоту перескоков атома\ispr{В исходном тексте
(с.~166) во втором множителе стоит $\EvF$ вместо $\Evm$.}
\begin{equation}\label{eq:Gamma-at}
\Gamma_{\text{ат}}=z\nu_0\Ex{-\frac{\EvF+\Evm}{kT}} .
\end{equation}

\paragraph{Коэффициент самодиффузии}
Пометим часть атомов (например, радиоактивным изотопом) и пусть их
концентрация меняется вдоль оси $x$. Возникнет диффузионный поток~— число
меченых атомов, проходящих за секунду через единичную площадку. По
\term{первому закону Фика} поток пропорционален градиенту концентрации:
\begin{equation}\label{eq:Fick1}
J=-D\frac{\d c}{\d x},
\end{equation}
где $D$~— коэффициент самодиффузии ($[D]=\text{см}^2/\text{с}$). Выразим $D$
через характеристики решётки на простейшей модели
(рис.~\ref{fig:two-planes}).

\begin{figure}[ht]
\centering
\begin{tikzpicture}[scale=0.75]
  \foreach \y in {0,...,4}{
    \node[atom=5mm] at (0,\y) {};
    \node[atom=5mm] at (1.6,\y) {};}
  \node[atomB=5mm] at (0,3) {}; \node[atomB=5mm] at (0,1) {}; \node[atomB=5mm] at (0,0) {};
  \node[atomB=5mm] at (1.6,2) {};
  \node[vac=5mm] at (1.6,1) {};
  \draw[arr, fmRed] (0.3,1.15) -- (1.3,1.15);
  \draw[arr, fmBlue] (1.3,2.15) -- (0.3,2.15);
  \node[lbl] at (0,-0.9) {плоскость 1};
  \node[lbl] at (1.6,-0.9) {плоскость 2};
  \node[lbl, align=left, right] at (2.6,3.2) {$c_1$~— концентрация меченых\\атомов в плоскости 1};
  \node[lbl, align=left, right] at (2.6,1.6) {$c_2=c_1+a\,\dfrac{\d c}{\d x}$};
  \draw[dimline] (0,4.6) -- node[above, slbl]{$a$} (1.6,4.6);
  \draw[arr] (-0.8,-1.6) -- node[below, slbl]{$x$} (2.4,-1.6);
\end{tikzpicture}
\caption{К выводу коэффициента диффузии. Меченые атомы (оранжевые)
перескакивают между соседними атомными плоскостями только через
вакансии}
\label{fig:two-planes}
\end{figure}

\begin{shagi}
\item Возьмём простую кубическую решётку с периодом $a$ и две соседние
плоскости $\{100\}$, перпендикулярные оси $x$. На единице площади
плоскости $1/a^2$ узлов. Доли меченых атомов в плоскостях $C_1$ и $C_2$.
\item У каждого атома плоскости~1 ровно один сосед в плоскости~2. Атом
перескакивает к нему, если этот узел пуст (вероятность
$\exp(-\EvF/kT)$) и если удалась попытка преодолеть барьер (частота
$\nu_0\exp(-\Evm/kT)$). Обозначим
$\Gamma=\nu_0\exp\bigl(-(\EvF+\Evm)/kT\bigr)$.
\item Поток меченых атомов слева направо $\Gamma\,C_1/a^2$, справа налево
$\Gamma\,C_2/a^2$. Результирующий поток
\[
J=\frac{\Gamma}{a^2}(C_1-C_2).
\]
\item Перейдём к объёмной концентрации $c=C/a^3$: градиент
$\d c/\d x=(C_2-C_1)/a^4$. Тогда
\[
J=-a^2\Gamma\,\frac{\d c}{\d x}\quad\Longrightarrow\quad D=a^2\Gamma .
\]
\end{shagi}

\begin{keyeq}\label{eq:Dsd}
D_{\text{сд}}=D_0\Ex{-\frac{\Ecd}{kT}},\qquad \Ecd=\EvF+\Evm,\qquad D_0\approx a^2\nu_0 .
\end{keyeq}
Энергия активации самодиффузии $\Ecd$ равна \emph{сумме} энергий
образования и миграции вакансии. Для меди $\Ecd\approx 1{,}28+0{,}7\approx
2{,}0\eV$, что хорошо согласуется с измеренным значением $\approx2{,}1\eV$.

\begin{podrobnee}[Общая формула и корреляционный множитель]
Для решётки любого кубического типа $D=\frac16 f\,\Gamma_{\text{ат}}\lambda^2$,
где $\Gamma_{\text{ат}}$~— полная частота перескоков~\eqref{eq:Gamma-at},
$\lambda$~— длина перескока. Для простой кубической решётки $z=6$,
$\lambda=a$, и мы получаем $D=a^2\Gamma$. Для ГЦК: $z=12$,
$\lambda=a/\sqrt2$, и снова $D=f a^2\nu_0\exp(-\Ecd/kT)$.

\emph{Корреляционный множитель} $f\le 1$ возникает потому, что перескоки
атома не независимы. Перескочив вперёд, атом оставляет вакансию прямо
позади себя, и пока она не ушла, следующий перескок с большой
вероятностью вернёт атом назад. Кроме того, атом, не успевший отдать
энергию соседям, может «отразиться» обратно. Для ГЦК-решётки
$f\approx 0{,}78$, для ОЦК $f\approx 0{,}73$; порядок величины $D$ не
меняется. В реальном $D_0$ входит также множитель
$\exp\bigl((S_v+S_m)/k\bigr)$ с энтропиями образования и миграции, поэтому
измеренные $D_0$ обычно в десятки раз больше, чем $a^2\nu_0$.
\end{podrobnee}

\begin{vnimanie}[Частая ошибка в толковании]
Формулу~\eqref{eq:Dsd} нередко объясняют так: «вакансия должна сначала
\emph{родиться} рядом с атомом (вероятность $e^{-\EvF/kT}$), а потом
\emph{переместиться} (вероятность $e^{-\Evm/kT}$)». Это неверно.
Вакансия не может спонтанно родиться в узле внутри кристалла: этому
мешает закон сохранения вещества (а рождение пары «вакансия +
межузельный атом» требует энергии, намного большей $\Ecd$). При
стационарной диффузии число вакансий постоянно, и множитель
$\exp(-\EvF/kT)$~— это просто \emph{вероятность застать} блуждающую
вакансию в соседнем узле.
\end{vnimanie}

\paragraph{Диффузия самих вакансий}
Коэффициент диффузии вакансий находится так же, но вакансии не нужно
ждать «партнёра»: любой соседний атом может с ней поменяться. Поэтому в
$D_v$ нет множителя $\exp(-\EvF/kT)$:
\begin{equation}\label{eq:Dv}
D_v=a^2\nu_0\Ex{-\frac{\Evm}{kT}},\qquad
\frac{D_{\text{сд}}}{D_v}=C_v .
\end{equation}
\begin{smysl}
Вакансия «бегает» в $1/C_v$ раз быстрее атома (у точки плавления~—
в $\sim10^4$ раз, при более низкой температуре~— во много больше).
Атом же почти всё время ждёт, пока к нему подойдёт вакансия.
\end{smysl}

\paragraph{Диффузия атомов замещения}
Атомы замещения, как и собственные атомы, сидят в узлах, и для них
основным тоже является вакансионный механизм:
\begin{equation}\label{eq:Di}
D_i=a^2\nu_0\Ex{-\frac{\EvF+E_i^{\mathrm m}}{kT}},
\end{equation}
где $E_i^{\mathrm m}$~— энергия активации перескока атома замещения в
соседнюю вакансию\ispr{В формуле на с.~169 оригинала в показателе стоит
$\Evm$ вместо $E_i^{\mathrm m}$, а вместо $a^2$~— $b^2$.}. Крупному атому
перескочить труднее ($E_i^{\mathrm m}>\Evm$), мелкому~— легче
($E_i^{\mathrm m}<\Evm$). Обычно размеры атомов различаются не сильно, и
эти энергии одного порядка.

%=======================================================================
\subsubsection{Подвижность межузельных атомов}\label{par:int-mobility}

Межузельные атомы намного подвижнее вакансий. Причина в том, что
гантельная и объёмно-центрированная конфигурации близки и по геометрии,
и по энергии. Самый лёгкий путь движения~— последовательная смена
конфигураций: гантель $\to$ объёмно-центрированная $\to$ гантель $\to\dots$
(рис.~\ref{fig:int-relay}). Движение носит \term{эстафетный} характер:
межузельный атом выталкивает соседа из узла, тот становится новым
межузельным атомом, и т.\,д.

\begin{figure}[ht]
\centering
\begin{tikzpicture}[scale=0.62]
  \def\cell#1#2{%
    \begin{scope}[shift={(#1,0)}]
    \foreach \x/\y in {0/0,2/0,0/2,2/2,1/1,3/1,-1/1,1/3,1/-1}{}%
    \foreach \x in {0,2}{\foreach \y in {0,2}{\node[atom=4mm] at (\x,\y) {};}}
    #2
    \end{scope}}
  % (а) межузельный атом в центре
  \cell{0}{\node[atom=4mm] at (1,-0.0) {}; \node[atom=4mm] at (0,1) {}; \node[atom=4mm] at (2,1) {};
           \node[atomI=4mm] at (1,1) {}; \node[atom=4mm] at (1,2) {};}
  % (б) гантель
  \cell{4.2}{\node[atom=4mm] at (1,-0.0) {}; \node[atom=4mm] at (0,1) {}; \node[atom=4mm] at (2,1) {};
           \node[vac=4mm] at (1,2) {}; \node[atomI=4mm] at (1,1.62) {}; \node[atomI=4mm] at (1,2.38) {};}
  % (в) снова центр — сдвинутый на полпериода
  \cell{8.4}{\node[atom=4mm] at (1,-0.0) {}; \node[atom=4mm] at (0,1) {}; \node[atom=4mm] at (2,1) {};
           \node[atom=4mm] at (1,2) {}; \node[atomI=4mm] at (1,3) {};}
  \node[lbl] at (1,-1.2) {а}; \node[lbl] at (5.2,-1.2) {б}; \node[lbl] at (9.4,-1.2) {в};
  \draw[arr, fmGray] (2.6,1) -- (3.6,1);
  \draw[arr, fmGray] (6.8,1) -- (7.8,1);
\end{tikzpicture}
\caption{Эстафетное движение межузельного атома (проекция на плоскость
$\{100\}$ ГЦК-решётки): объёмно-центрированная конфигурация~(а)
переходит в гантельную~(б), затем снова в объёмно-центрированную,
сдвинутую на полпериода~(в). Красным выделены атомы, принадлежащие
дефекту}
\label{fig:int-relay}
\end{figure}

Если гантель выгоднее объёмно-центрированной конфигурации, то на кривой
энергии межузельного атома в середине пути появляется
\emph{промежуточный минимум} (рис.~\ref{fig:int-profile}). Барьер,
который нужно преодолеть, оказывается гораздо ниже, чем для вакансии:
\begin{glavnoe}
Энергия миграции межузельного атома очень мала: $\Eim\approx 0{,}1\eV$
(против $\Evm\approx 1\eV$). Поэтому межузельные атомы начинают
двигаться уже при очень низких температурах (в меди~— при
$\sim 40$~К).
\end{glavnoe}

\begin{figure}[ht]
\centering
\begin{tikzpicture}
\begin{axis}[width=0.66\textwidth, height=6.4cm,
  xmin=-0.05, xmax=2.05, ymin=0, ymax=3.3, axis lines=left,
  xtick={0,0.5,1,1.5,2}, xticklabels={$0$,$\tfrac12a$,$a$,$\tfrac32a$,$2a$},
  ytick={0.6,2.2,2.55}, yticklabels={$\EvF$,$E^{\mathrm F}_{i1}$,$E^{\mathrm F}_{i2}$},
  xlabel={положение дефекта $x$}, ylabel={$E$}]
  \addplot[fmRed, very thick, domain=0:2, samples=200]
    {2.2+0.55*sin(deg(2*pi*x))^2+0.35*sin(deg(pi*x))^2};
  \addplot[fmBlue, very thick, domain=0:2, samples=200]
    {0.6+0.9*sin(deg(pi*x))^2};
  \draw[dimline] (axis cs:1.95,2.2) -- node[left, font=\small]{$\Eim$} (axis cs:1.95,2.86);
  \draw[dimline] (axis cs:1.5,0.6) -- node[right, font=\small]{$\Evm$} (axis cs:1.5,1.5);
  \node[fmRed, font=\small] at (axis cs:1,3.1) {межузельный атом};
  \node[fmBlue, font=\small] at (axis cs:1,0.35) {вакансия};
\end{axis}
\end{tikzpicture}
\caption{Энергия межузельного атома (вверху) и вакансии (внизу) при их
движении вдоль $x$. $E^{\mathrm F}_{i1}$ и $E^{\mathrm F}_{i2}$~—
энергии образования межузельного атома в гантельной и
объёмно-центрированной конфигурации. Промежуточный минимум делает барьер
для межузельного атома низким}
\label{fig:int-profile}
\end{figure}

\paragraph{Краудион}
Краудион занимает промежуточное положение между точечным и линейным
дефектом: смещения вокруг него велики вдоль ряда (до $\sim10\,a$) и малы
поперёк. Движется он тоже необычно: вдоль своего ряда очень легко
(барьер $\sim0{,}1\Evm$, его можно преодолеть даже за счёт кинетической
энергии), а в других направлениях~— практически никак. Это пример
\emph{одномерной} миграции.

\begin{itogi}
\begin{itemize}
\item $\EvF\approx E_c-W$: простая модель разрыва связей даёт $E_c$,
релаксация снижает значение в 2,5--3 раза. Практическое правило:
$\EvF\approx10\,k\Tpl\sim1\eV$. Для межузельного атома
$\EiF\approx3$--$5\eV$.
\item Равновесная концентрация: $C=\exp(S/k)\exp(-E^{\mathrm F}/kT)$.
У точки плавления $C_v\sim10^{-4}$; равновесных межузельных атомов
практически нет.
\item В сплаве вакансии могут концентрироваться у примесных атомов, если
$B_{iv}/kT>\ln(1/zC_i)$.
\item Частота перескоков вакансии $z\nu_0\exp(-\Evm/kT)$; атома~—
$z\nu_0\exp\bigl(-(\EvF+\Evm)/kT\bigr)$.
\item Самодиффузия: $D=D_0\exp(-\Ecd/kT)$, $\Ecd=\EvF+\Evm$; вакансии
диффундируют в $1/C_v$ раз быстрее атомов.
\item Межузельные атомы очень подвижны: $\Eim\sim0{,}1\eV$ (эстафетный
механизм); краудион движется одномерно.
\end{itemize}
\end{itogi}

\begin{proverka}
\item Почему оценка $\EvF\approx E_c$ по разрыву связей завышена?
\item Во сколько раз изменится равновесная концентрация вакансий в
меди при нагреве от 1000 до 1300~К?
\item Почему в равновесии межузельных атомов практически нет, а вакансии
есть?
\item Какой процесс медленнее: перескок вакансии в соседний узел или
перескок данного атома? Во сколько раз?
\item Почему энергия активации самодиффузии равна сумме $\EvF+\Evm$, а не
одному $\Evm$?
\item Почему межузельные атомы подвижнее вакансий?
\end{proverka}
